\chapter*{Abstract}

This thesis studies adversarial robustness, privacy, and reproducibility in safety-critical machine learning systems, with particular emphasis on computer vision, anomaly detection, and evasion attacks through a series of papers.
The work begins by analysing the practical costs and benefits of defence strategies against adversarial attacks, revealing that common robustness metrics are poor indicators of real-world adversarial performance (Paper I). 
Through large-scale experiments, it further demonstrates that adversarial examples can often be generated in linear time, granting attackers a computational advantage over defenders (Paper II).
To address this, a novel metric—the Training Rate and Survival Heuristic (TRASH)—was developed to predict model failure under attack and facilitate early rejection of vulnerable architectures (Paper III).
This metric was then extended to real-world cost, showing that adversarial robustness can be improved using low-cost, low-precision hardware without sacrificing accuracy (Paper IV).

Beyond robustness, the thesis tackles privacy by designing a lightweight, client-side spam detection model that preserves user data and resists several classes of attacks without requiring server-side computation (Paper V).
Recognizing the need for reproducible and auditable experiments in safety-critical contexts, the thesis also presents \texttt{deckard}, a declarative software framework for distributed and robust machine learning experimentation (Paper VI).

Together, these contributions offer empirical techniques for evaluating and improving model robustness, propose a privacy-preserving classification strategy, and deliver practical tooling for reproducible experimentation. 
Ultimately, this thesis advances the goal of building machine learning systems that are not only accurate, but also robust, reproducible, and trustworthy.

\chapter*{Sammanfattning}

Denna avhandling studerar robusthet, integritet och reproducerbarhet i säker\-hetskritisk maskininlärning, med särskild tonvikt på datorseende, avvikelse\-detektering och undvikande attacker.

Arbetet inleds med att analysera de praktiska kostnaderna och fördelarna med försvarsstrategier mot attacker, vilket visar att vanliga mått på robusthet är dåliga indikatorer på verklig prestanda i attacker (Artikel~I).
Genom storskaliga experiment visar arbetet vidare att exempel på attacker ofta kan genereras i linjär tid, vilket ger angripare en beräkningsfördel gentemot försvar\-are (Artikel~II).
För att hantera detta presenterar avhandlingen ett nytt mått – Training Rate and Survival Heuristic (TRASH) – för att förutsäga modellfel under attack och underlätta tidigt avvisande av sårbara arkitekturer (Artikel~III).
Detta mått utvidgades sedan till verkliga kostnader, vilket visar att robusthet i attacker kan förbättras med hjälp av billig hårdvara med låg precision utan att offra noggrannheten (Artikel~IV).

Utöver robusthet behandlar avhandlingen integritet genom att utforma en lättviktig klientbaserad modell för spamdetektering som bevarar användardata och står emot flera klasser av attacker utan att kräva att beräkningar görs på serversidan (Artikel~V).
Som svar på behovet av reproducerbara och gransk\-ningsbara experiment i säkerhetskritiska sammanhang presenterar avhandlingen även \texttt{deckard}, ett deklarativt programvaruramverk för distribuerade och robusta maskininlärningsexperiment (Artikel~VI).

Tillsammans erbjuder dessa bidrag empiriska tekniker för att utvärdera och förbättra modellers robusthet, föreslår en integritetsbevarande klassificeringsstrategi och levererar praktiska verktyg för reproducerbara experiment.
Sammantaget främjar avhandlingen målet att bygga maskininlärningssystem som inte bara är korrekta, utan också robusta, reproducerbara och pålitliga.